\addtocontents{toc}{\protect\setcounter{tocdepth}{0}}
\chapter[Anexo: Diccionario de Datos]{ANEXO: Diccionario de Datos del Modelo de Datos}

Este anexo describe las tablas de la base de datos PostgreSQL 15 definidas en el
script de esquema del proyecto. Los nombres de tablas y columnas son los que
usa el servidor. Las columnas marcadas como cifradas contienen texto cifrado con
AES-256-GCM en el dispositivo del cuidador. El servidor no puede descifrarlas y
rechaza la escritura de datos clínicos en claro.

Las tablas pacientes, registros\_clinicos, recordatorios y
conversaciones guardan sus datos sensibles en una columna cifrados de
tipo jsonb, que reúne claves cuyo nombre termina en \_cifrado
o \_cifrada. Para cada una de esas tablas se describen primero las columnas y
después las claves de su columna cifrados. Las columnas en claro son las que
el servidor necesita para relacionar, ordenar y filtrar. Los tipos de fecha
usan zona horaria (timestamptz) y los identificadores de usuario son uuid.

\section{Resumen de las tablas}

\begingroup
\small
\tablazebra
\begin{xltabular}{\textwidth}{P{0.90}P{1.10}}
    \caption{Tablas de la base de datos y su entidad de dependencia.}
    \label{tab:dicc-jerarquia}\\
    \toprule
    \filaencabezado \encab{Tabla} & \encab{Depende de} \\
    \midrule
    \endfirsthead
    \rowcolor{white}\multicolumn{2}{@{}l}{\footnotesize\itshape Tabla \thetable{} (continuación)} \\
    \toprule
    \filaencabezado \encab{Tabla} & \encab{Depende de} \\
    \midrule
    \endhead
    \midrule
    \rowcolor{white}\multicolumn{2}{r@{}}{\footnotesize\itshape Continúa en la página siguiente} \\
    \endfoot
    \bottomrule
    \endlastfoot
    usuarios & Raíz: el cuidador \\
    pacientes & usuarios (cuidador) \\
    registros\_\allowbreak{}clinicos & pacientes \\
    recordatorios & pacientes (programados por el cuidador) \\
    conversaciones & usuarios (cuidador) \\
    favoritos & usuarios y material\_educativo (tabla puente) \\
    material\_\allowbreak{}educativo & Transversal, compartida y de solo lectura \\
    claves\_\allowbreak{}datos & usuarios (una por cuidador) \\
    tokens\_\allowbreak{}restablecer & usuarios \\
    tokens\_\allowbreak{}correo & usuarios \\
    esperas\_\allowbreak{}recuperacion & Independiente (limita reenvíos) \\
\end{xltabular}
\endgroup

Las preguntas frecuentes que usan el chat y la pantalla de preguntas frecuentes
no residen en la base de datos. Se incluyen en la propia aplicación, en la lista
preguntasFrecuentes del módulo faq, y el chat elige la
respuesta por coincidencia de palabras clave. Tampoco existe una tabla de
síntomas, porque el catálogo ESAS forma parte de la aplicación.

\section{Cuidador (usuarios)}

\begingroup
\small
\tablazebra
\begin{xltabular}{\textwidth}{P{1.05}P{0.55}P{0.32}P{1.58}P{0.50}}
    \caption{Diccionario de datos de la tabla usuarios.}
    \label{tab:dicc-usuarios}\\
    \toprule
    \filaencabezado \encab{Columna} & \encab{Tipo} & \encab{Nulo} & \encab{Descripción} & \encab{Cifrado} \\
    \midrule
    \endfirsthead
    \rowcolor{white}\multicolumn{5}{@{}l}{\footnotesize\itshape Tabla \thetable{} (continuación)} \\
    \toprule
    \filaencabezado \encab{Columna} & \encab{Tipo} & \encab{Nulo} & \encab{Descripción} & \encab{Cifrado} \\
    \midrule
    \endhead
    \midrule
    \rowcolor{white}\multicolumn{5}{r@{}}{\footnotesize\itshape Continúa en la página siguiente} \\
    \endfoot
    \bottomrule
    \endlastfoot
    id & uuid & No & Identificador del cuidador. Clave primaria, generada por el servidor. & No \\
    correo & text & No & Correo principal de acceso. Es único. & No \\
    contrasena\_\allowbreak{}hash & text & Sí & Contraseña almacenada como hash bcrypt. & No \\
    creado\_\allowbreak{}en & timestamptz & Sí & Fecha de creación de la cuenta. & No \\
    version\_\allowbreak{}encriptacion & int & Sí & Versión del esquema de cifrado aplicado (2 por defecto). & No \\
    cifrados & jsonb & Sí & Texto cifrado de los datos personales del cuidador (claves de la tabla siguiente). & Sí \\
    correo\_\allowbreak{}verificado & boolean & Sí & Indica si el correo principal fue confirmado mediante enlace. & No \\
    correo\_\allowbreak{}respaldo\_\allowbreak{}hash & text & Sí & HMAC del correo de respaldo confirmado. Nunca se guarda el correo en claro. & No \\
    pendiente\_\allowbreak{}correo & text & Sí & Correo nuevo a la espera de confirmación. & No \\
    pendiente\_\allowbreak{}correo\_\allowbreak{}tipo & text & Sí & Tipo de cambio pendiente (cambio de correo principal o respaldo). & No \\
    pendiente\_\allowbreak{}correo\_\allowbreak{}solicitado\_\allowbreak{}en & timestamptz & Sí & Fecha de la solicitud pendiente. & No \\
    respaldo\_\allowbreak{}pendiente\_\allowbreak{}servidor & boolean & Sí & Marca que existe un respaldo por confirmar. & No \\
\end{xltabular}
\endgroup

\begingroup
\small
\tablazebra
\begin{xltabular}{\textwidth}{P{1.05}P{0.55}P{0.32}P{1.58}P{0.50}}
    \caption{Claves de la columna cifrados de usuarios.}
    \label{tab:dicc-usuarios-cifrados}\\
    \toprule
    \filaencabezado \encab{Columna} & \encab{Tipo} & \encab{Nulo} & \encab{Descripción} & \encab{Cifrado} \\
    \midrule
    \endfirsthead
    \rowcolor{white}\multicolumn{5}{@{}l}{\footnotesize\itshape Tabla \thetable{} (continuación)} \\
    \toprule
    \filaencabezado \encab{Columna} & \encab{Tipo} & \encab{Nulo} & \encab{Descripción} & \encab{Cifrado} \\
    \midrule
    \endhead
    \midrule
    \rowcolor{white}\multicolumn{5}{r@{}}{\footnotesize\itshape Continúa en la página siguiente} \\
    \endfoot
    \bottomrule
    \endlastfoot
    nombre\_\allowbreak{}cifrado & texto cifrado & Sí & Nombre completo del cuidador. & Sí \\
    telefono\_\allowbreak{}cifrado & texto cifrado & Sí & Teléfono del cuidador. & Sí \\
    relacion\_\allowbreak{}cifrada & texto cifrado & Sí & Parentesco con el paciente (madre, padre, tutor u otro). & Sí \\
    direccion\_\allowbreak{}cifrada & texto cifrado & Sí & Dirección del cuidador. & Sí \\
    correo\_\allowbreak{}respaldo\_\allowbreak{}cifrado & texto cifrado & Sí & Correo de respaldo, cifrado para que el cliente pueda mostrarlo. & Sí \\
\end{xltabular}
\endgroup

\section{Paciente (pacientes)}

\begingroup
\small
\tablazebra
\begin{xltabular}{\textwidth}{P{1.05}P{0.55}P{0.32}P{1.58}P{0.50}}
    \caption{Diccionario de datos de la tabla pacientes.}
    \label{tab:dicc-pacientes}\\
    \toprule
    \filaencabezado \encab{Columna} & \encab{Tipo} & \encab{Nulo} & \encab{Descripción} & \encab{Cifrado} \\
    \midrule
    \endfirsthead
    \rowcolor{white}\multicolumn{5}{@{}l}{\footnotesize\itshape Tabla \thetable{} (continuación)} \\
    \toprule
    \filaencabezado \encab{Columna} & \encab{Tipo} & \encab{Nulo} & \encab{Descripción} & \encab{Cifrado} \\
    \midrule
    \endhead
    \midrule
    \rowcolor{white}\multicolumn{5}{r@{}}{\footnotesize\itshape Continúa en la página siguiente} \\
    \endfoot
    \bottomrule
    \endlastfoot
    usuario\_\allowbreak{}id & uuid & No & Cuidador propietario. Clave foránea a usuarios, con eliminación en cascada. Parte de la clave primaria. & No \\
    id & text & No & Identificador del paciente. Parte de la clave primaria. & No \\
    notificaciones\_\allowbreak{}activas & boolean & Sí & Indica si las notificaciones del paciente están habilitadas (verdadero por defecto). & No \\
    maximo\_\allowbreak{}registros\_\allowbreak{}dia & int & Sí & Límite de registros clínicos programados por día (3 por defecto). & No \\
    archivado & boolean & Sí & Marca el paciente como archivado y lo excluye de las listas activas. & No \\
    creado\_\allowbreak{}en & timestamptz & Sí & Fecha de alta del paciente. & No \\
    version\_\allowbreak{}encriptacion & int & Sí & Versión del esquema de cifrado. & No \\
    cifrados & jsonb & Sí & Texto cifrado de los datos del paciente (claves de la tabla siguiente). & Sí \\
\end{xltabular}
\endgroup

\begingroup
\small
\tablazebra
\begin{xltabular}{\textwidth}{P{1.05}P{0.55}P{0.32}P{1.58}P{0.50}}
    \caption{Claves de la columna cifrados de pacientes.}
    \label{tab:dicc-pacientes-cifrados}\\
    \toprule
    \filaencabezado \encab{Columna} & \encab{Tipo} & \encab{Nulo} & \encab{Descripción} & \encab{Cifrado} \\
    \midrule
    \endfirsthead
    \rowcolor{white}\multicolumn{5}{@{}l}{\footnotesize\itshape Tabla \thetable{} (continuación)} \\
    \toprule
    \filaencabezado \encab{Columna} & \encab{Tipo} & \encab{Nulo} & \encab{Descripción} & \encab{Cifrado} \\
    \midrule
    \endhead
    \midrule
    \rowcolor{white}\multicolumn{5}{r@{}}{\footnotesize\itshape Continúa en la página siguiente} \\
    \endfoot
    \bottomrule
    \endlastfoot
    nombre\_\allowbreak{}cifrado & texto cifrado & Sí & Nombre del paciente. & Sí \\
    rut\_\allowbreak{}cifrado & texto cifrado & Sí & RUT del paciente. & Sí \\
    edad\_\allowbreak{}cifrada & texto cifrado & Sí & Edad del paciente. & Sí \\
    diagnostico\_\allowbreak{}cifrado & texto cifrado & Sí & Diagnóstico oncológico. & Sí \\
    fase\_\allowbreak{}tratamiento\_\allowbreak{}cifrado & texto cifrado & Sí & Fase del tratamiento. & Sí \\
    centro\_\allowbreak{}salud\_\allowbreak{}nombre\_\allowbreak{}cifrado & texto cifrado & Sí & Nombre del centro de salud. & Sí \\
    centro\_\allowbreak{}salud\_\allowbreak{}direccion\_\allowbreak{}cifrado & texto cifrado & Sí & Dirección del centro de salud. & Sí \\
    centro\_\allowbreak{}salud\_\allowbreak{}telefono\_\allowbreak{}cifrado & texto cifrado & Sí & Teléfono del centro de salud. & Sí \\
    contacto\_\allowbreak{}emergencia\_\allowbreak{}nombre\_\allowbreak{}cifrado & texto cifrado & Sí & Nombre del contacto de emergencia. & Sí \\
    contacto\_\allowbreak{}emergencia\_\allowbreak{}telefono\_\allowbreak{}cifrado & texto cifrado & Sí & Teléfono del contacto de emergencia. & Sí \\
\end{xltabular}
\endgroup

\section{Registro clínico (registros\_clinicos)}

\begingroup
\small
\tablazebra
\begin{xltabular}{\textwidth}{P{1.05}P{0.55}P{0.32}P{1.58}P{0.50}}
    \caption{Diccionario de datos de la tabla registros\_clinicos.}
    \label{tab:dicc-registros}\\
    \toprule
    \filaencabezado \encab{Columna} & \encab{Tipo} & \encab{Nulo} & \encab{Descripción} & \encab{Cifrado} \\
    \midrule
    \endfirsthead
    \rowcolor{white}\multicolumn{5}{@{}l}{\footnotesize\itshape Tabla \thetable{} (continuación)} \\
    \toprule
    \filaencabezado \encab{Columna} & \encab{Tipo} & \encab{Nulo} & \encab{Descripción} & \encab{Cifrado} \\
    \midrule
    \endhead
    \midrule
    \rowcolor{white}\multicolumn{5}{r@{}}{\footnotesize\itshape Continúa en la página siguiente} \\
    \endfoot
    \bottomrule
    \endlastfoot
    usuario\_\allowbreak{}id & uuid & No & Cuidador propietario. Parte de la clave primaria y de la clave foránea compuesta a pacientes. & No \\
    paciente\_\allowbreak{}id & text & No & Paciente al que pertenece el registro. Parte de la clave primaria y de la clave foránea a pacientes, con eliminación en cascada. & No \\
    id & text & No & Identificador del registro. Parte de la clave primaria. & No \\
    fecha & timestamptz & Sí & Fecha de la valoración. Se usa para filtrar por rango. & No \\
    creado\_\allowbreak{}en & timestamptz & Sí & Fecha y hora de creación. Ordena el historial. & No \\
    tipo\_\allowbreak{}registro & text & Sí & Tipo de registro clínico (programado o extra). & No \\
    nivel\_\allowbreak{}alerta & text & Sí & Nivel calculado por el motor de reglas clínicas: normal, alerta o crítico. & No \\
    version\_\allowbreak{}encriptacion & int & Sí & Versión del esquema de cifrado. & No \\
    cifrados & jsonb & Sí & Texto cifrado del contenido clínico (claves de la tabla siguiente). & Sí \\
\end{xltabular}
\endgroup

\begingroup
\small
\tablazebra
\begin{xltabular}{\textwidth}{P{1.05}P{0.55}P{0.32}P{1.58}P{0.50}}
    \caption{Claves de la columna cifrados de registros\_clinicos.}
    \label{tab:dicc-registros-cifrados}\\
    \toprule
    \filaencabezado \encab{Columna} & \encab{Tipo} & \encab{Nulo} & \encab{Descripción} & \encab{Cifrado} \\
    \midrule
    \endfirsthead
    \rowcolor{white}\multicolumn{5}{@{}l}{\footnotesize\itshape Tabla \thetable{} (continuación)} \\
    \toprule
    \filaencabezado \encab{Columna} & \encab{Tipo} & \encab{Nulo} & \encab{Descripción} & \encab{Cifrado} \\
    \midrule
    \endhead
    \midrule
    \rowcolor{white}\multicolumn{5}{r@{}}{\footnotesize\itshape Continúa en la página siguiente} \\
    \endfoot
    \bottomrule
    \endlastfoot
    signos\_\allowbreak{}vitales\_\allowbreak{}cifrado & texto cifrado & Sí & Temperatura, frecuencia cardíaca, saturación y frecuencia respiratoria, serializados como JSON antes de cifrar. & Sí \\
    sintomas\_\allowbreak{}cifrado & texto cifrado & Sí & Lista de síntomas ESAS con intensidad 0--10 y nota por síntoma, serializada y cifrada. & Sí \\
    contenido\_\allowbreak{}registro\_\allowbreak{}cifrado & texto cifrado & Sí & Observaciones del registro. & Sí \\
    mensaje\_\allowbreak{}alerta\_\allowbreak{}cifrado & texto cifrado & Sí & Mensaje de alerta asociado al nivel detectado. & Sí \\
\end{xltabular}
\endgroup

Los signos vitales y los síntomas se guardan en una única clave cifrada cada uno, porque el cliente los serializa como JSON antes de cifrar. Un registro sin signos vitales omite esa clave.

\section{Recordatorio (recordatorios)}

\begingroup
\small
\tablazebra
\begin{xltabular}{\textwidth}{P{1.05}P{0.55}P{0.32}P{1.58}P{0.50}}
    \caption{Diccionario de datos de la tabla recordatorios.}
    \label{tab:dicc-recordatorios}\\
    \toprule
    \filaencabezado \encab{Columna} & \encab{Tipo} & \encab{Nulo} & \encab{Descripción} & \encab{Cifrado} \\
    \midrule
    \endfirsthead
    \rowcolor{white}\multicolumn{5}{@{}l}{\footnotesize\itshape Tabla \thetable{} (continuación)} \\
    \toprule
    \filaencabezado \encab{Columna} & \encab{Tipo} & \encab{Nulo} & \encab{Descripción} & \encab{Cifrado} \\
    \midrule
    \endhead
    \midrule
    \rowcolor{white}\multicolumn{5}{r@{}}{\footnotesize\itshape Continúa en la página siguiente} \\
    \endfoot
    \bottomrule
    \endlastfoot
    usuario\_\allowbreak{}id & uuid & No & Cuidador que programa el recordatorio. Parte de la clave primaria y de la clave foránea compuesta a pacientes. & No \\
    paciente\_\allowbreak{}id & text & No & Paciente al que se dirige el recordatorio. Clave foránea a pacientes, con eliminación en cascada. & No \\
    id & text & No & Identificador del recordatorio. Parte de la clave primaria. & No \\
    activo & boolean & Sí & Indica si el recordatorio está activo (verdadero por defecto). & No \\
    creado\_\allowbreak{}en & timestamptz & Sí & Fecha de creación. & No \\
    version\_\allowbreak{}encriptacion & int & Sí & Versión del esquema de cifrado. & No \\
    cifrados & jsonb & Sí & Texto cifrado del recordatorio (claves de la tabla siguiente). & Sí \\
\end{xltabular}
\endgroup

\begingroup
\small
\tablazebra
\begin{xltabular}{\textwidth}{P{1.05}P{0.55}P{0.32}P{1.58}P{0.50}}
    \caption{Claves de la columna cifrados de recordatorios.}
    \label{tab:dicc-recordatorios-cifrados}\\
    \toprule
    \filaencabezado \encab{Columna} & \encab{Tipo} & \encab{Nulo} & \encab{Descripción} & \encab{Cifrado} \\
    \midrule
    \endfirsthead
    \rowcolor{white}\multicolumn{5}{@{}l}{\footnotesize\itshape Tabla \thetable{} (continuación)} \\
    \toprule
    \filaencabezado \encab{Columna} & \encab{Tipo} & \encab{Nulo} & \encab{Descripción} & \encab{Cifrado} \\
    \midrule
    \endhead
    \midrule
    \rowcolor{white}\multicolumn{5}{r@{}}{\footnotesize\itshape Continúa en la página siguiente} \\
    \endfoot
    \bottomrule
    \endlastfoot
    titulo\_\allowbreak{}cifrado & texto cifrado & Sí & Título del recordatorio. & Sí \\
    descripcion\_\allowbreak{}cifrada & texto cifrado & Sí & Descripción del recordatorio. & Sí \\
    datos\_\allowbreak{}cifrados & texto cifrado & Sí & Programación serializada como JSON antes de cifrar: tipo (medicamento, medición, cita u otro), fecha y hora, días de repetición, recurrencia y a quién se dirige. & Sí \\
\end{xltabular}
\endgroup

\section{Conversación (conversaciones)}
\label{sec:conversaciones}

\begingroup
\small
\tablazebra
\begin{xltabular}{\textwidth}{P{1.05}P{0.55}P{0.32}P{1.58}P{0.50}}
    \caption{Diccionario de datos de la tabla conversaciones.}
    \label{tab:dicc-conversaciones}\\
    \toprule
    \filaencabezado \encab{Columna} & \encab{Tipo} & \encab{Nulo} & \encab{Descripción} & \encab{Cifrado} \\
    \midrule
    \endfirsthead
    \rowcolor{white}\multicolumn{5}{@{}l}{\footnotesize\itshape Tabla \thetable{} (continuación)} \\
    \toprule
    \filaencabezado \encab{Columna} & \encab{Tipo} & \encab{Nulo} & \encab{Descripción} & \encab{Cifrado} \\
    \midrule
    \endhead
    \midrule
    \rowcolor{white}\multicolumn{5}{r@{}}{\footnotesize\itshape Continúa en la página siguiente} \\
    \endfoot
    \bottomrule
    \endlastfoot
    usuario\_\allowbreak{}id & uuid & No & Cuidador propietario. Clave foránea a usuarios, con eliminación en cascada. Parte de la clave primaria. & No \\
    id & text & No & Identificador de la conversación. Parte de la clave primaria. & No \\
    ultima\_\allowbreak{}actividad & timestamptz & Sí & Fecha de la última actividad. Ordena las conversaciones. & No \\
    creado\_\allowbreak{}en & timestamptz & Sí & Fecha de creación. & No \\
    version\_\allowbreak{}encriptacion & int & Sí & Versión del esquema de cifrado. & No \\
    cifrados & jsonb & Sí & Texto cifrado de la conversación (claves de la tabla siguiente). & Sí \\
\end{xltabular}
\endgroup

\begingroup
\small
\tablazebra
\begin{xltabular}{\textwidth}{P{1.05}P{0.55}P{0.32}P{1.58}P{0.50}}
    \caption{Claves de la columna cifrados de conversaciones.}
    \label{tab:dicc-conversaciones-cifrados}\\
    \toprule
    \filaencabezado \encab{Columna} & \encab{Tipo} & \encab{Nulo} & \encab{Descripción} & \encab{Cifrado} \\
    \midrule
    \endfirsthead
    \rowcolor{white}\multicolumn{5}{@{}l}{\footnotesize\itshape Tabla \thetable{} (continuación)} \\
    \toprule
    \filaencabezado \encab{Columna} & \encab{Tipo} & \encab{Nulo} & \encab{Descripción} & \encab{Cifrado} \\
    \midrule
    \endhead
    \midrule
    \rowcolor{white}\multicolumn{5}{r@{}}{\footnotesize\itshape Continúa en la página siguiente} \\
    \endfoot
    \bottomrule
    \endlastfoot
    titulo\_\allowbreak{}cifrado & texto cifrado & Sí & Título de la conversación. & Sí \\
    mensajes\_\allowbreak{}cifrado & texto cifrado & Sí & Mensajes con su texto e indicador de emisor, serializados como JSON antes de cifrar. Su orden es la posición en la lista. & Sí \\
\end{xltabular}
\endgroup

A diferencia de los registros y los recordatorios, las conversaciones dependen del cuidador y no del paciente, porque el chat es general para la cuenta, igual que el contenido educativo.

\section{Material educativo (material\_educativo)}

\begingroup
\small
\tablazebra
\begin{xltabular}{\textwidth}{P{1.05}P{0.55}P{0.32}P{1.58}P{0.50}}
    \caption{Diccionario de datos de la tabla material\_educativo (transversal, de solo lectura desde la aplicación).}
    \label{tab:dicc-material}\\
    \toprule
    \filaencabezado \encab{Columna} & \encab{Tipo} & \encab{Nulo} & \encab{Descripción} & \encab{Cifrado} \\
    \midrule
    \endfirsthead
    \rowcolor{white}\multicolumn{5}{@{}l}{\footnotesize\itshape Tabla \thetable{} (continuación)} \\
    \toprule
    \filaencabezado \encab{Columna} & \encab{Tipo} & \encab{Nulo} & \encab{Descripción} & \encab{Cifrado} \\
    \midrule
    \endhead
    \midrule
    \rowcolor{white}\multicolumn{5}{r@{}}{\footnotesize\itshape Continúa en la página siguiente} \\
    \endfoot
    \bottomrule
    \endlastfoot
    id & text & No & Identificador del material. Clave primaria. & No \\
    titulo & text & Sí & Título del material. & No \\
    categoria & text & Sí & Categoría (filtro de la biblioteca). & No \\
    tema & text & Sí & Tema o etiqueta del material. & No \\
    cuerpo & text & Sí & Cuerpo del contenido, con secciones tituladas. & No \\
    url\_\allowbreak{}imagen & text & Sí & Ruta de la imagen destacada. & No \\
    url\_\allowbreak{}archivo & text & Sí & Ruta del archivo adjunto (video, guía, infografía o documento), servido por la API. & No \\
    url\_\allowbreak{}miniatura & text & Sí & Ruta de la miniatura. & No \\
    tipo\_\allowbreak{}archivo & text & Sí & Tipo del archivo adjunto. & No \\
    tamano\_\allowbreak{}bytes & bigint & Sí & Tamaño del archivo adjunto. & No \\
    creado\_\allowbreak{}en & timestamptz & Sí & Fecha de publicación. & No \\
\end{xltabular}
\endgroup

La aplicación no escribe en esta tabla. La API la carga al iniciar, a partir de un
catálogo incluido en el repositorio del proyecto, solo si está vacía, y expone
los archivos adjuntos desde su propio almacenamiento. Las rutas se guardan
relativas y la API las entrega como direcciones absolutas.

\section{Favoritos (favoritos)}

\begingroup
\small
\tablazebra
\begin{xltabular}{\textwidth}{P{1.05}P{0.55}P{0.32}P{1.58}P{0.50}}
    \caption{Diccionario de datos de la tabla favoritos (tabla puente entre cuidador y material educativo).}
    \label{tab:dicc-favoritos}\\
    \toprule
    \filaencabezado \encab{Columna} & \encab{Tipo} & \encab{Nulo} & \encab{Descripción} & \encab{Cifrado} \\
    \midrule
    \endfirsthead
    \rowcolor{white}\multicolumn{5}{@{}l}{\footnotesize\itshape Tabla \thetable{} (continuación)} \\
    \toprule
    \filaencabezado \encab{Columna} & \encab{Tipo} & \encab{Nulo} & \encab{Descripción} & \encab{Cifrado} \\
    \midrule
    \endhead
    \midrule
    \rowcolor{white}\multicolumn{5}{r@{}}{\footnotesize\itshape Continúa en la página siguiente} \\
    \endfoot
    \bottomrule
    \endlastfoot
    usuario\_\allowbreak{}id & uuid & No & Cuidador que marcó el favorito. Clave foránea a usuarios, con eliminación en cascada. Parte de la clave primaria. & No \\
    material\_\allowbreak{}id & text & No & Material marcado. Referencia lógica a material\_educativo. Parte de la clave primaria. & No \\
    creado\_\allowbreak{}en & timestamptz & Sí & Fecha en que se marcó. & No \\
\end{xltabular}
\endgroup

\section{Tablas de seguridad de la cuenta}

Estas tablas no forman parte del modelo conceptual, porque el cliente nunca las
lee ni las escribe de forma directa. La administra únicamente la API, que las
usa para entregar la clave de datos, restablecer contraseñas y confirmar correos.

\begingroup
\small
\tablazebra
\begin{xltabular}{\textwidth}{P{1.05}P{0.55}P{0.32}P{1.58}P{0.50}}
    \caption{Diccionario de datos de la tabla claves\_datos.}
    \label{tab:dicc-claves}\\
    \toprule
    \filaencabezado \encab{Columna} & \encab{Tipo} & \encab{Nulo} & \encab{Descripción} & \encab{Cifrado} \\
    \midrule
    \endfirsthead
    \rowcolor{white}\multicolumn{5}{@{}l}{\footnotesize\itshape Tabla \thetable{} (continuación)} \\
    \toprule
    \filaencabezado \encab{Columna} & \encab{Tipo} & \encab{Nulo} & \encab{Descripción} & \encab{Cifrado} \\
    \midrule
    \endhead
    \midrule
    \rowcolor{white}\multicolumn{5}{r@{}}{\footnotesize\itshape Continúa en la página siguiente} \\
    \endfoot
    \bottomrule
    \endlastfoot
    usuario\_\allowbreak{}id & uuid & No & Cuidador dueño de la clave. Clave primaria y clave foránea a usuarios, con eliminación en cascada. & No \\
    clave\_\allowbreak{}sellada & text & No & Clave de datos del cliente (32 bytes) sellada con AES-256-GCM mediante la clave maestra del servidor, con formato iv.etiqueta.cifrado en base64. & Sí \\
    creado\_\allowbreak{}en & timestamptz & Sí & Fecha de creación. & No \\
\end{xltabular}
\endgroup

\begingroup
\small
\tablazebra
\begin{xltabular}{\textwidth}{P{1.05}P{0.55}P{0.32}P{1.58}P{0.50}}
    \caption{Diccionario de datos de la tabla tokens\_restablecer.}
    \label{tab:dicc-restablecer}\\
    \toprule
    \filaencabezado \encab{Columna} & \encab{Tipo} & \encab{Nulo} & \encab{Descripción} & \encab{Cifrado} \\
    \midrule
    \endfirsthead
    \rowcolor{white}\multicolumn{5}{@{}l}{\footnotesize\itshape Tabla \thetable{} (continuación)} \\
    \toprule
    \filaencabezado \encab{Columna} & \encab{Tipo} & \encab{Nulo} & \encab{Descripción} & \encab{Cifrado} \\
    \midrule
    \endhead
    \midrule
    \rowcolor{white}\multicolumn{5}{r@{}}{\footnotesize\itshape Continúa en la página siguiente} \\
    \endfoot
    \bottomrule
    \endlastfoot
    token\_\allowbreak{}hash & text & No & SHA-256 del token enviado por correo. Clave primaria. El token en claro nunca se guarda. & No \\
    usuario\_\allowbreak{}id & uuid & No & Cuidador que solicitó el restablecimiento. Clave foránea a usuarios, con eliminación en cascada. & No \\
    expira\_\allowbreak{}en & timestamptz & No & Vencimiento del enlace (una hora desde su emisión). & No \\
    usado & boolean & Sí & Indica si el enlace ya se utilizó (de un solo uso). & No \\
\end{xltabular}
\endgroup

\begingroup
\small
\tablazebra
\begin{xltabular}{\textwidth}{P{1.05}P{0.55}P{0.32}P{1.58}P{0.50}}
    \caption{Diccionario de datos de la tabla tokens\_correo.}
    \label{tab:dicc-tcorreo}\\
    \toprule
    \filaencabezado \encab{Columna} & \encab{Tipo} & \encab{Nulo} & \encab{Descripción} & \encab{Cifrado} \\
    \midrule
    \endfirsthead
    \rowcolor{white}\multicolumn{5}{@{}l}{\footnotesize\itshape Tabla \thetable{} (continuación)} \\
    \toprule
    \filaencabezado \encab{Columna} & \encab{Tipo} & \encab{Nulo} & \encab{Descripción} & \encab{Cifrado} \\
    \midrule
    \endhead
    \midrule
    \rowcolor{white}\multicolumn{5}{r@{}}{\footnotesize\itshape Continúa en la página siguiente} \\
    \endfoot
    \bottomrule
    \endlastfoot
    token\_\allowbreak{}hash & text & No & SHA-256 del token enviado por correo. Clave primaria. & No \\
    usuario\_\allowbreak{}id & uuid & No & Cuidador al que corresponde el enlace. Clave foránea a usuarios, con eliminación en cascada. & No \\
    tipo & text & No & Tipo de confirmación: verificacion (registro), cambio (correo principal) o respaldo. & No \\
    correo\_\allowbreak{}nuevo & text & Sí & Correo que se activará al confirmar (solo para el cambio del correo principal). & No \\
    respaldo\_\allowbreak{}hash & text & Sí & HMAC del correo de respaldo que se activará al confirmar. & No \\
    creado\_\allowbreak{}en & timestamptz & Sí & Fecha de emisión. & No \\
    expira\_\allowbreak{}en & timestamptz & No & Vencimiento del enlace (una hora desde su emisión). & No \\
    usado & boolean & Sí & Indica si el enlace ya se utilizó (de un solo uso). & No \\
\end{xltabular}
\endgroup

\begingroup
\small
\tablazebra
\begin{xltabular}{\textwidth}{P{1.05}P{0.55}P{0.32}P{1.58}P{0.50}}
    \caption{Diccionario de datos de la tabla esperas\_recuperacion.}
    \label{tab:dicc-esperas}\\
    \toprule
    \filaencabezado \encab{Columna} & \encab{Tipo} & \encab{Nulo} & \encab{Descripción} & \encab{Cifrado} \\
    \midrule
    \endfirsthead
    \rowcolor{white}\multicolumn{5}{@{}l}{\footnotesize\itshape Tabla \thetable{} (continuación)} \\
    \toprule
    \filaencabezado \encab{Columna} & \encab{Tipo} & \encab{Nulo} & \encab{Descripción} & \encab{Cifrado} \\
    \midrule
    \endhead
    \midrule
    \rowcolor{white}\multicolumn{5}{r@{}}{\footnotesize\itshape Continúa en la página siguiente} \\
    \endfoot
    \bottomrule
    \endlastfoot
    hash & text & No & Hash del identificador de la solicitud de recuperación. Clave primaria. & No \\
    ultimo\_\allowbreak{}envio\_\allowbreak{}en & timestamptz & Sí & Fecha del último envío. Limita la frecuencia de reenvíos. & No \\
\end{xltabular}
\endgroup

La clave de datos de cada cuidador se genera la primera vez que la solicita la
aplicación (\texttt{POST /clave-datos}) y se guarda sellada con la clave maestra
definida en el entorno del servidor. Si esa clave maestra se pierde, las claves
selladas no pueden abrirse. La aplicación no realiza intercambio de claves
propio: recibe la clave de datos por una conexión autenticada y la conserva en el
almacenamiento seguro del teléfono.

\addtocontents{toc}{\protect\setcounter{tocdepth}{2}}
